% ===========================================================================
%  Diagrama a bloques - Practica 4: Convertidores de nivel y motor con PC817
%  Programa ESP32-S3 (practica4_convertidores_pc817.ino)
%  Optoelectronica (V3736) - Dr. Ruben Estrada Marmolejo, CUCEI-UDG
% ===========================================================================
\documentclass[border=14pt]{standalone}
\usepackage[T1]{fontenc}
\usepackage[utf8]{inputenc}
\usepackage[spanish,es-noshorthands]{babel}
\usepackage{amsmath}
\usepackage{tikz}
\usetikzlibrary{positioning, arrows.meta, shapes.geometric, calc, fit, backgrounds}

% ---------- Paleta ----------
\definecolor{colEsp}{HTML}{1F6FB2}       % ESP32 / firmware (azul)
\definecolor{colEspBg}{HTML}{E3F0FA}
\definecolor{colOpto}{HTML}{B5651D}      % barrera PC817 (ocre)
\definecolor{colOptoBg}{HTML}{FBEEDD}
\definecolor{colC1}{HTML}{0E8A6A}        % circuito 1 (verde)
\definecolor{colC1Bg}{HTML}{E1F5EE}
\definecolor{colC2}{HTML}{6A3FB5}        % circuito 2 (morado)
\definecolor{colC2Bg}{HTML}{EEE6F8}
\definecolor{colC3}{HTML}{B32651}        % circuito 3 motor (magenta)
\definecolor{colC3Bg}{HTML}{FBE3EC}
\definecolor{colSetup}{HTML}{555F6B}     % setup (gris)
\definecolor{colSetupBg}{HTML}{EDEFF2}
\definecolor{colLinea}{HTML}{2B2F36}

\begin{document}
\begin{tikzpicture}[
  font=\sffamily,
  >=Latex,
  every node/.style={transform shape},
  base/.style={rounded corners=3pt, draw, line width=0.7pt, align=center,
               inner sep=4pt, minimum height=11mm, text width=30mm},
  esp/.style={base, draw=colEsp, fill=colEspBg, text=colLinea},
  setup/.style={base, draw=colSetup, fill=colSetupBg, text=colLinea, text width=40mm},
  opto/.style={base, draw=colOpto, fill=colOptoBg, text=colLinea, text width=26mm},
  c1/.style={base, draw=colC1, fill=colC1Bg, text=colLinea},
  c2/.style={base, draw=colC2, fill=colC2Bg, text=colLinea},
  c3/.style={base, draw=colC3, fill=colC3Bg, text=colLinea, text width=26mm},
  flecha/.style={->, line width=0.9pt, draw=colLinea},
  etiq/.style={font=\sffamily\scriptsize\itshape, text=colLinea, fill=white, inner sep=1pt}
]

% ----- Columnas (x) y carriles (y) -----
\def\xCtrl{0}      % GPIO del ESP32
\def\xOpto{5.4}    % barrera PC817
\def\xPow{10.8}    % cargas / campo
\def\xPowB{15.6}   % segunda etapa (motor)
\def\yA{0}         % carril C1
\def\yB{-2.6}      % carril C3 (motor)
\def\yC{-5.2}      % carril C2

% ================= FIRMWARE (cabecera) =================
\node[setup] (loop) at (\xCtrl,2.7)
   {\textbf{setup() + loop()}\\ {\scriptsize Serial 115200 \textbullet\ ledcAttach 1\,kHz}\\ {\scriptsize despachador: c1 / c2 / vel / stop}};

% ================= CARRIL A: C1 (3.3 -> 24 V) =================
\node[esp] (g1) at (\xCtrl,\yA) {\textbf{GPIO4 OUT}\\ {\scriptsize digitalWrite}};
\node[opto] (o1) at (\xOpto,\yA) {\textbf{PC817 \#1}\\ {\scriptsize R\textsubscript{in}=220\,$\Omega$}};
\node[c1] (load1) at (\xPow,\yA) {\textbf{Salida 24 V}\\ {\scriptsize R\textsubscript{C}=10\,k a +24\,V}\\ {\scriptsize nivel invertido}};

% ================= CARRIL B: C3 (motor PWM) =================
\node[esp] (g3) at (\xCtrl,\yB) {\textbf{GPIO6 PWM}\\ {\scriptsize ledcWrite 0..1023}\\ {\scriptsize compensa inv.}};
\node[opto] (o3) at (\xOpto,\yB) {\textbf{PC817 \#3}\\ {\scriptsize R\textsubscript{in}=220\,$\Omega$}\\ {\scriptsize aisla el PWM}};
\node[c3] (mosfet) at (\xPow,\yB) {\textbf{MOSFET + flyback}\\ {\scriptsize R\textsubscript{C}=10\,k a +5\,V}\\ {\scriptsize compuerta}};
\node[c3] (motor) at (\xPowB,\yB) {\textbf{Motor DC 5 V}\\ {\scriptsize motorreductor TT}\\ {\scriptsize vel $\propto$ duty}};

% ================= CARRIL C: C2 (24 -> 3.3 V, fluye a la izquierda) =================
\node[esp] (g2) at (\xCtrl,\yC) {\textbf{GPIO5 IN}\\ {\scriptsize digitalRead + niega}};
\node[opto] (o2) at (\xOpto,\yC) {\textbf{PC817 \#2}\\ {\scriptsize R\textsubscript{in}=4.7\,k$\Omega$}};
\node[c2] (src24) at (\xPow,\yC) {\textbf{Senal 24 V}\\ {\scriptsize entrada de campo}\\ {\scriptsize excita el LED}};

% ===================== FLECHAS =====================
% firmware -> GPIOs
\draw[flecha] (loop.south) -- ++(0,-0.5) -| (g1.north);
\draw[flecha] (loop.south) -- ++(0,-0.5) -| (g3.north);
% A: ->
\draw[flecha] (g1) -- (o1) node[etiq,pos=0.5]{LED};
\draw[flecha] (o1) -- (load1) node[etiq,pos=0.5]{24\,V};
% B: ->
\draw[flecha] (g3) -- (o3) node[etiq,pos=0.5]{PWM};
\draw[flecha] (o3) -- (mosfet);
\draw[flecha] (mosfet) -- (motor);
% C: <- (entrada que viaja hacia el ESP)
\draw[flecha] (src24) -- (o2) node[etiq,pos=0.5]{LED};
\draw[flecha] (o2) -- (g2) node[etiq,pos=0.5]{3.3\,V};

% ===================== BANDAS VERTICALES (zonas) =====================
\begin{scope}[on background layer]
  % lado control
  \node[rounded corners=4pt, fill=colEspBg!45, draw=colEsp!35, dashed,
        fit=(loop)(g1)(g3)(g2), inner sep=4mm,
        label={[colEsp,font=\sffamily\bfseries\small]below:LADO CONTROL (ESP32-S3, GND\textsubscript{ESP})}] (zCtrl) {};
  % barrera
  \node[rounded corners=4pt, fill=colOptoBg!40, draw=colOpto!45, dashed,
        fit=(o1)(o3)(o2), inner sep=4mm,
        label={[colOpto,font=\sffamily\bfseries\small]below:BARRERA\\ (PC817)}] (zOpto) {};
  % lado potencia / campo
  \node[rounded corners=4pt, fill=colC3Bg!32, draw=colC3!35, dashed,
        fit=(load1)(mosfet)(motor)(src24), inner sep=4mm,
        label={[colC3,font=\sffamily\bfseries\small]below:LADO POTENCIA / CAMPO (GND\textsubscript{24V}, GND\textsubscript{motor})}] (zPow) {};
\end{scope}

% ===================== TITULO =====================
\node[font=\sffamily\bfseries\large, text=colLinea, anchor=south west]
   at ([yshift=5mm]zCtrl.north west)
   {Programa ESP32-S3 -- Pr\'actica 4: convertidores de nivel 3.3\,$\leftrightarrow$\,24\,V y motor PWM con PC817};

\end{tikzpicture}
\end{document}
